\input format.tex

\usepackage{graphicx}
\graphicspath{{cores/}}

\usepackage{environ}
\usepackage{colortbl,array,booktabs}
\usepackage{tabularx}

\colorlet{TablaBordeSuperior}{topcolor}
\colorlet{TablaBordeInferior}{topcolor}
\colorlet{TablaCentroSuperior}{blue!1}
\colorlet{TablaCentroInferior}{blue!20}
\colorlet{FuenteCabeceraTabla}{white}

\newcolumntype{M}[1]{>{\centering\arraybackslash}m{#1}}

\tcbset{rtab/.style={
freelance,
frame code={
 \path[top color=topcolor,bottom color=topcolor]
   ([yshift=-#1*(\baselineskip+2pt)]interior.north west) --
   ([yshift=-#1*(\baselineskip+2pt)]interior.north east) {[sharp corners]--
    ([yshift=3pt]interior.north east) --
    ([yshift=3pt]interior.north west)} -- cycle;

  },
interior code={},
 }
}

\newcommand\fuentecabecera[1]{\textcolor{black}{\textbf{#1}}}

\NewEnviron{RCtable}[1]
  {%
    \begin{tcolorbox}[#1]\BODY\end{tcolorbox}%
  }



\begin{document}

\vspace*{5mm}
%% 各章节
\setlength{\arrayrulewidth}{1pt}
\fontsize{9.3pt}{11pt}\selectfont
\color{gray2}

{\noindent\bf\sanhao 二、代谢分析及免疫}

\vspace*{4mm}

\begin{spacing}{1.5}
\begin{LRaside}[.8]{2.1 消化和吸收总评}
\noindent
\includegraphics[width=\linewidth]{hanliang.pdf}
\asidebreak %
蛋白质、脂类、碳水化合物是维持生命及健康的重要营养物质。蛋白质是生命的物质基础，构成人体各个组织器官，实现各种生理功能；脂类是人体主要储能物质；碳水化合物是机体最重要的直接供能物质。肠道菌群在蛋白质、脂类、碳水化合物的消化、代谢过程中起着重要的作用。
\end{LRaside}
\end{spacing}

%\vspace*{4mm}
%\noindent
%2.1 您的常见肠道菌含量
%\vspace*{1.5mm}

%\begin{RCtable}{
%  rtab=1,
%  tabularx*={\arrayrulecolor{gray2}}%
%    {M{0.3\linewidth} M{0.31\linewidth}M{0.34\linewidth}}
%}
%\fuentecabecera{测试项目} & \fuentecabecera{健康状态提示} & \fuentecabecera{检测结果评价} \\[3pt]
%\bf{测试项目} & \bf{健康状态提示} & \bf{检测结果评价} \\[3pt]
%\end{RCtable}

\begin{tctabularx}{tabularx={m{4.8cm}<{\centering}m{5.2cm}<{\centering}m{4.55cm}<{\centering}}}
&&
\\[-6pt]
\cellcolor{topcolor} \raisebox{2.614pt}{\color{gray0}\bf 检测项目} &
\cellcolor{topcolor} \raisebox{2.614pt}{\color{gray0}\bf 健康状态提示} &
\cellcolor{topcolor} \raisebox{2.614pt}{\color{gray0}\bf 检测结果评价}
\end{tctabularx}

\vspace*{-4.25mm}
\fontsize{8.8pt}{11pt}\selectfont
\arrayrulecolor{gray2}
\begin{longtable}{m{4.8cm}m{5.2cm}<{\centering}m{0cm}@{}m{4.61cm}<{\centering}}
\hline
\parbox[c]{\hsize}{\vskip7pt {\lantxh 蛋白降解产物\\肠道菌群氨基酸代谢能力} \vskip7pt} & \parbox[c]{\hsize}{\vskip7pt\centerline{\raisebox{-1.5ex}{\includegraphics[width=5cm,height=0.55cm]{bar.pdf}}}\vskip7pt}  &
\hspace*{-4.83cm}\raisebox{-0.45ex}{\includegraphics{smile.pdf}}
 & \begin{minipage}{4.60cm}\begin{center}{{\color{green}\lantxh 偏高{\\ \bahao 有利于人体获得必需氨基酸}} }\end{center} \end{minipage} \\
% & \begin{minipage}{4.60cm}\begin{center}{{\lantxh 偏高{\\ \bahao 有利于人体获得必需氨基酸}} }\end{center} \end{minipage} \\
% & \begin{minipage}{4.60cm}\begin{center}{{\color{green}\lantxh 偏高 \\ \bahao 有利于人体获得必需氨基酸} }\end{center} \end{minipage} \\
\hline
\parbox[c]{\hsize}{\vskip7pt {\lantxh 脂肪降解产物\\肠道菌群脂肪类相关代谢能力} \vskip7pt} & \parbox[c]{\hsize}{\vskip7pt\centerline{\raisebox{-1.5ex}{\includegraphics[width=5cm,height=0.55cm]{bar.pdf}}}\vskip7pt}  &
\hspace*{-4.83cm}\raisebox{-0.45ex}{\includegraphics{smile.pdf}}
 & \begin{minipage}{4.60cm}\begin{center}{{\color{green}\lantxh 偏高{\\ \bahao 有利于人体获得必需脂肪酸}} }\end{center} \end{minipage} \\
% & \begin{minipage}{4.60cm}\begin{center}{{\lantxh 偏高{\\ \bahao 有利于人体获得必需脂肪酸}} }\end{center} \end{minipage} \\
% & \begin{minipage}{4.60cm}\begin{center}{{\color{green}\lantxh 偏高 \\ \bahao 有利于人体获得必需脂肪酸} }\end{center} \end{minipage} \\
\hline
\parbox[c]{\hsize}{\vskip7pt {\lantxh 碳水化合物降解产物\\肠道菌群纤维素等碳水化合物代谢能力} \vskip7pt} & \parbox[c]{\hsize}{\vskip7pt\centerline{\raisebox{-1.5ex}{\includegraphics[width=5cm,height=0.55cm]{bar.pdf}}}\vskip7pt}  &
\hspace*{-4.83cm}\raisebox{-0.45ex}{\includegraphics{smile.pdf}}
 & \begin{minipage}{4.60cm}\begin{center}{{\color{green}\lantxh 偏高{\\ \bahao 有利于维持肠道稳态和菌群平衡}} }\end{center} \end{minipage} \\
% & \begin{minipage}{4.60cm}\begin{center}{{\lantxh 偏高{\\ \bahao 有利于维持肠道稳态和菌群平衡}} }\end{center} \end{minipage} \\
% & \begin{minipage}{4.60cm}\begin{center}{{\color{green}\lantxh 偏高 \\ \bahao 有利于维持肠道稳态和菌群平衡} }\end{center} \end{minipage} \\
\hline
\multicolumn{3}{l}{}\\
\multicolumn{3}{l}{\includegraphics[scale=1]{greencircle.pdf} {\fontsize{9.3pt}{11pt}\selectfont\bf 绿色表示健康}   \includegraphics[scale=1]{orangecircle.pdf}  {\fontsize{9.3pt}{11pt}\selectfont\bf 橙色表示需要关注}  \includegraphics[scale=1]{redcircle.pdf} {\fontsize{9.3pt}{11pt}\selectfont\bf 红色表示有风险}}
\end{longtable}

\vspace*{6mm}
\begin{spacing}{1.5}
\begin{LRaside}[.8]{结果分析}
\noindent
\includegraphics[width=\linewidth]{result_fenbu.pdf}
\asidebreak %
综合您的检测结果，
您肠道内的蛋白降解产物、脂肪降解产物、碳水化合物降解产物指标优于参考人群，有利于获得必需氨基酸、获得必需脂肪酸、肠道稳态等。有利于肠道和人体健康，请您继续保持良好的饮食生活习惯。
\end{LRaside}
\end{spacing}

\end{document}

